\section[Stage 040]{Stage 040: Generalized Selected-Branch Normalization with Source/Loading Mismatch}
\label{stage:040}

\stagefield{Part / anchor}{Part III; primary anchor \resultanchor{MTDC-T6}.}
\claimstatus{\StatusExactClosure{} for the rank-one selected-mode algebra with distinct source and loading directions.}
\stagefield{Purpose}{Stage~040 replaces the flat one-vector selected branch by a source/loading mismatch law.  It derives deformed branch functions \(F_{q,\eta_s}\) and \(G_q\), then specializes them to the split--\(U\) continuum through \(R_U\).}
\stagefield{Inputs}{The selected lower eigenvalue \(\lambda_-=A_0(1-\xi)\), loading vector \(z=z_0(1,q)^T\), and source vector \(s=s_0(1,t)^T\).}

\stagefield{Verification}{SymPy audit: \StageFile{scripts/moving_throat_pde_stage040_generalized_selected_branch_sympy_audit.py}; transcript: \StageFile{scripts/output/moving_throat_pde_stage040_generalized_selected_branch_sympy_audit.txt}.  Mathematica audit: \StageFile{mathematica/moving_throat_pde_stage040_generalized_selected_branch_mathematica_audit.wl}; transcript: \StageFile{mathematica/output/moving_throat_pde_stage040_generalized_selected_branch_mathematica_audit.txt}.}

\subsection*{Derivation}

For a single rank-one loading, the required loading is
\begin{equation}
\label{eq:app-stage040-alpha-req}
\alpha_{\rm req}=\frac{A_0\xi(\xi+\delta)}{z_0^2[\delta+(1+q^2)\xi]},
\qquad
0\leq\xi<1.
\end{equation}
The selected eigenvector has ratio
\begin{equation}
\label{eq:app-stage040-evec}
\frac{e_1}{e_0}=\frac{q\xi}{\delta+\xi}.
\end{equation}
With \(\eta_s=s_1z_1/(s_0z_0)\), the generalized normalization and baseline loading functions are
\begin{equation}
\label{eq:app-stage040-Fqeta}
\boxed{
F_{q,\eta_s}(\xi,\delta)
=\frac{[\delta+(1+q^2)\xi]^2[\delta+(1+\eta_s)\xi]^2}
{(1-\xi)[(\delta+\xi)^2+q^2\xi^2]^2}},
\end{equation}
\begin{equation}
\label{eq:app-stage040-Gq}
\boxed{
G_q(\xi,\delta)=\frac{\xi(\xi+\delta)}{\delta+(1+q^2)\xi}}.
\end{equation}
For the split--\(U\) continuum, using \(\lambda_0=\kappa_1^2/\kappa_0^2=2/9\),
\begin{equation}
\label{eq:app-stage040-qeta-RU}
q=-\frac{\sqrt2}{3}R_U,
\qquad
\eta_s=\frac{2}{9}R_U.
\end{equation}
Therefore the one-parameter split--\(U\) functions are
\begin{equation}
\label{eq:app-stage040-FU}
\boxed{
F_U(\xi,\delta;R_U)=
\frac{[9\delta+(9+2R_U^2)\xi]^2[9\delta+(9+2R_U)\xi]^2}
{81(1-\xi)[9\delta^2+18\delta\xi+(9+2R_U^2)\xi^2]^2}},
\end{equation}
\begin{equation}
\label{eq:app-stage040-GU}
\boxed{
G_U(\xi,\delta;R_U)=\frac{9\xi(\xi+\delta)}{9\delta+(9+2R_U^2)\xi}}.
\end{equation}
The flat branch is recovered at \(R_U=1\).

\stagefield{Output}{The generalized selected functions \eqref{eq:app-stage040-Fqeta}--\eqref{eq:app-stage040-Gq} and their split--\(U\) specialization \eqref{eq:app-stage040-FU}--\eqref{eq:app-stage040-GU}.}
\stagefield{Downstream use}{Stages~041--044 add a second support direction and determine the actual softening depth.}

